\chapter{Conclusion}
%%%%%%%%%%%%%%%%%%%%

\section{Summary}

This thesis set out to answer whether per-user variation in analog mouse-click
telemetry is structured enough to support automatic labeling, statistical modeling,
and firmware parameter optimization beyond static default settings.  The answer is
yes on all three fronts, confirmed by a complete pipeline evaluated on a corpus of
132 annotated recordings from 39 participants, collected across two data-collection
campaigns: an in-house controlled study and a public LAN gaming event.

\paragraph{Labeling pipeline.}
An algorithmic bootstrapping heuristic combined with iterative model-assisted
annotation produced 39,088 verified left-button and 6,645 right-button clicks
at 1\,kHz.  The three-phase workflow (heuristic bootstrap $\to$ model-assisted
labeling $\to$ human review) substantially reduced annotation burden relative to
labeling from scratch while maintaining biomechanical fidelity: the heuristic alone
achieves a pooled deployment F1 of only 19.9\,\% across the full corpus, showing that
human review is essential for the firmware-unregistered and noise-confounded cases
that are most informative for model training.

\paragraph{Click-intent classification.}
Trained on the resulting labeled corpus, the Random Forest achieves a deployment
F1 of 97.0\,\% on the left channel, with make- and break-boundary bias well
within the $\pm4$\,ms matching radius used to score correctness --- though this
reflects tight agreement with the ground-truth boundary annotation itself
rather than an independently established perceptual threshold for click
timing. The RF modestly outperforms Gradient Boosted Trees on this single
held-out split despite the two achieving near-identical grouped
cross-validation generalization (Tables~\ref{tab:deployment}
and~\ref{tab:cv}), and trains roughly twice as fast as the FCN on the left
channel --- a real but, as Section~\ref{sec:eval-discussion} notes, secondary
consideration next to the cross-validated generalization gap --- together
favoring the RF as the default architecture for deployment.
The dilated FCN, despite a 3,060-frame receptive field, provides no measurable
gain on the left channel in the single-test comparison and collapses under
grouped cross-validation on that same channel (F1 $59.7 \pm 17.4$ percentage
points (pp)); this is consistent with large-context models requiring
proportionally large annotated corpora, though an underexplored hyperparameter
space means a poorly tuned model cannot be ruled out
(Section~\ref{sec:conclusion-lim}).

\paragraph{Firmware threshold optimizer.}
Using the classifier's click-boundary predictions, the exhaustive FSM simulator
reduced optimizer loss by 49\,\% (mean), increased
the TP rate from 90.9\,\% to 97.5\,\%, and reduced
mean make latency from 27.0\,ms to 15.4\,ms --- an 11.6\,ms reduction that, as
Section~\ref{sec:eval-discussion} discusses, is best read as a measured
latency improvement rather than a confirmed perceptible gain, absent a
citable perceptual threshold for this context, and whose persistence beyond
the session it was fit on has not yet been tested (Section~\ref{sec:conclusion-lim}).

\paragraph{End-to-end user study.}
A double-blind study with 26 participants (13 calibrated, 13 placebo) tested
whether these optimizer-predicted gains are perceptible in gameplay. The
calibrated group
showed a statistically significant, large improvement in clicks-per-second
($+10.7\,\%$ vs.\ $-2.5\,\%$ for placebo, $p=0.035$, $d=0.92$), directionally
favorable but inconclusive results on a tracking task, and a non-significant
accuracy trade-off on a fast-flicking task accompanied by a higher misclick count.
The study is consistent with the optimizer behaving well end-to-end, but does not
by itself confirm it. It produced a non-degenerate spread of per-user
actuation-point recommendations, though part of that spread is inflated by a
CPS-specific artifact: participants who pre-press before the round starts leave
the optimizer only one low-information calibration point to fit, pushing
recommendations toward the depth extremes rather than reflecting genuine
per-user signal. Recommended reactive-trigger sensitivity, by contrast, did not
differentiate participants at all --- 25 of 26 were assigned the lowest
available level --- which may reflect a genuinely uniform optimum across users
or limited sensitivity of the underlying signal at this stage of the pipeline
(Section~\ref{sec:eval-user-study}). The study also surfaced a
concrete risk (click-timing gains trading off against flick accuracy) that the
offline evaluation alone could not have revealed.

\section{Limitations}
\label{sec:conclusion-lim}

The limitations below fall into four groups: data and labeling quality (heuristic
unreliability, right-channel imbalance, single-hardware coverage, annotation
ambiguity), modeling (FCN hyperparameter search, absence of session-level context,
the binary click-matching criterion), the optimizer's untested cross-session
robustness, and the end-to-end user study's statistical power and validity
(gaming-experience moderation, sample size, participant overlap).

\paragraph{Heuristic unreliability across participants.}
Section~\ref{sec:eval-discussion} shows that the heuristic bootstrap's per-participant
deployment F1 spans 0.1\,\% to 75.0\,\% on the left channel and 0.0\,\% to 43.5\,\% on
the right, and that this variance is not explained by any collected demographic or
self-reported attribute. The practical consequence is that the
heuristic cannot serve as a standalone labeling or deployment strategy for any given
user: its reliability cannot be predicted in advance from the attributes we
collected. This is a limitation of the heuristic itself, not merely a design
choice made around it: the three-phase annotation pipeline already treats the
heuristic as a bootstrap to be corrected by human review and a learned classifier
rather than as a deployable detector in its own right, and this finding is what
would block relaxing that correction step in the future, not just what motivated
it originally.

\paragraph{Right-channel data imbalance.}
The right-button corpus is approximately six times smaller than the left (6,645 vs.\
39,088 clicks), causing the FCN and the heuristic to degrade significantly on that
channel.  A balanced collection protocol --- or a dedicated right-click minigame that
naturally elicits more right-button activity --- is needed for practical
right-channel deployment.

\paragraph{Single hardware platform.}
All data was collected on a single HITS-enabled PRO~X2 Superstrike mouse.
Generalization to other analog input devices with different sensor resolutions,
quantization levels, or firmware architectures has not been tested.  The
feature-engineering pipeline (120-level depth, derivative channels) is designed to
be device-agnostic, but empirical transfer must be verified.

\paragraph{Annotation ambiguity at onset boundaries.}
The $\pm 4$\,ms onset tolerance absorbs natural annotator variability, but for
borderline clicks at the edge of the biomechanical window the correct label boundary
is genuinely ambiguous.  The asymmetric soft-label targets implemented for the GBT
model (Section~\ref{sec:modeling}) could partially address this, but were not used
for the reported results, and a formal inter-annotator agreement study was not
conducted; the reported metrics therefore include an irreducible component of
label noise.

\paragraph{FCN hyperparameter underexploration.}
The dilated residual FCN was trained for only 16--17 epochs with a fixed
hyperparameter set due to the cost of training (\textasciitilde44 minutes per
run, already with RTX~5080 CUDA acceleration). No systematic search over
architecture depth, dropout rate, learning rate schedule, or number of dilation
stages was performed. The observed grouped cross-validation collapse
($59.7 \pm 17.4$\,pp) may therefore reflect a poorly tuned model rather than a
fundamental architectural limitation. An Optuna sweep over the FCN's
hyperparameter space, combined with a larger right-button corpus, is the most
likely path to closing the performance gap with the Random Forest; since GPU
acceleration is already in use, further speedup would have to come from a
smaller search space or distributed sweeps rather than additional hardware.

\paragraph{Absence of session-level temporal context.}
All three architectures operate on a local window spanning only a few tens of
milliseconds of hand-crafted features, augmented with $\pm$10\,ms context probes
(Section~\ref{sec:eval-channel-ablation}). The Random Forest has no access to
session-level dynamics: the current click's depth profile is evaluated in isolation
from the cumulative fatigue drift, the inter-click tempo, or the running statistics
of the session so far. Biomechanical research suggests that finger press depth and
velocity evolve systematically over extended play sessions, and that click intent
signatures cluster differently at the start versus end of a 40-minute session; a
classifier blind to this drift cannot distinguish a genuine late-session change in
intent from noise. Section~\ref{sec:conclusion-fw} outlines a concrete feature-level
extension addressing this gap.

\paragraph{Binary click-matching criterion.}
The dual-radius matching criterion used for deployment F1 (Section~\ref{sec:eval-setup})
classifies each (predicted, ground-truth) pair as a binary hit or miss based on whether
\emph{both} onsets fall within a fixed 4\,ms radius of their respective ground-truth
boundary centers. This treats a prediction whose make-onset misses the radius by 1\,ms
the same as one that misses by 50\,ms, and discards any information about near-misses
that fail only one of the two radius checks. A graded similarity metric --- such as a
distance-weighted match score or a duration-error penalty --- would provide a more
informative picture of where boundary errors concentrate and would allow distinguishing
a nearly correct prediction from a completely wrong one.

\paragraph{Optimizer cross-session robustness untested.}
The corpus-wide optimizer benchmark (Section~\ref{sec:eval-opt}) fits and scores
each configuration on the same session's data, which establishes the achievable
magnitude of the loss reduction and make-latency improvement but not whether a
configuration fit on one session continues to perform as well when applied
prospectively to a later session from the same participant. Since within-user
behavior may drift between sessions, this remains the primary open question
about how much of the reported 49\,\% loss reduction and 11.6\,ms make-latency
gain would hold up in a deployed, cross-session setting.

\paragraph{No moderating effect of gaming experience detected.}
A post-hoc mixed ANOVA (group $\times$ self-reported gamer status) found no
significant interaction on any task (Section~\ref{sec:eval-user-study}), meaning
the calibration effect could not be shown to differ between self-identified
gamers and non-gamers. Given the recruitment imbalance noted below and the small
sample size, this null result should be read as inconclusive rather than as
evidence that gaming experience does not moderate the effect.

\paragraph{Underpowered user study.}
The end-to-end user study (Section~\ref{sec:eval-user-study}) used 13 participants
per group, well below the $\sim$29 per group ($n=58$ total) that a Lehr's-formula
power calculation at the observed effect sizes ($d\approx0.75$) indicates is
needed for adequately powered ($1-\beta=0.80$) comparisons. Consistent with this,
the 95\,\% confidence intervals on Cohen's $d$ for all five non-CPS metrics
include zero (Table~\ref{tab:user-study}), meaning none of those effects is
statistically distinguishable from no effect at this sample size. Between-subject assignment left
uncontrolled variance in baseline skill level; no warmup round preceded the timed
tasks, so while the between-group $\Delta$ comparison controls for the average
learning effect between calibration and test rounds, it does not control for a
possible group $\times$ learning-effect interaction (e.g.\ if the calibrated
configuration itself changes how much a participant improves with practice);
and gamer/non-gamer counts were not balanced in
advance across groups. The Tracking and Flicking results should therefore be read
as hints rather than conclusive findings, and the Flicking accuracy/misclick
trade-off in particular warrants replication before it is treated as a genuine
effect of lowering the actuation point.

\paragraph{Participant overlap between training corpus and user study.}
Some of the 26 end-to-end user-study participants also contributed to the
39-participant corpus used to train the Random Forest classifier
(Section~\ref{sec:eval-user-study}), so the RF may have been exposed to some
study participants' own data during training. The exact overlap was not
tracked at the individual-participant level, so its magnitude cannot be
quantified after the fact, and no sensitivity analysis excluding overlapping
participants was run. This should therefore be kept in mind as a directional
caveat when interpreting the measured CPS gain: part of it may reflect
familiarity with those participants' click patterns rather than genuine
per-user generalization to unseen users.

\section{Future Work}
\label{sec:conclusion-fw}

\paragraph{Cross-device and cross-game generalization.}
Replicating the data collection on Hall-effect keyboards and alternative mouse
switches would test whether the trained feature representations transfer, and whether
a single universal model is feasible or per-device fine-tuning is required.  The
game-genre effect (FPS vs.\ MOBA click patterns, as observed in the PolyLAN dataset)
is also worth modeling explicitly to produce genre-aware recommendations.

\paragraph{On-device inference.}
The current pipeline is entirely offline.  Quantizing the Random Forest to INT8 and
deploying it on the peripheral MCU would enable real-time per-frame classification at
1\,kHz without a host connection, matching the latency requirements discussed in the
Background chapter and opening the path to adaptive threshold adjustment during play.

\paragraph{Adaptive recalibration.}
Neuromotor fatigue shifts the analog depth signal over the course of a gaming session
--- specifically, the resting drift channel rises as the flexor muscles tire.  A
lightweight online loop that monitors the running mean of the resting-drift channel
and nudges the actuation threshold accordingly could maintain the optimized
configuration as the user's biomechanical state evolves.

\paragraph{Session-aware global feature extraction.}
The most impactful architectural extension for the Random Forest would be to
append session-level summary features to each per-frame feature vector:
cumulative click count, running mean and variance of the resting-drift channel,
recent inter-click interval (a proxy for play tempo and arousal), and a fatigue
indicator derived from the slow drift of the resting depth baseline over the
session. These features would allow the model to adapt its decision boundary as
both the hardware signal and the user's neuromuscular state drift across a
40-minute session, addressing the absence of long-range temporal context that
the probe ablation identified as the primary bottleneck.

\paragraph{Two-stage cascade with boundary refinement.}
The current pipeline treats detection and boundary placement as a single
classification step. A two-stage architecture would separate these concerns:
Stage~1 uses the existing RF (fast, reliable at detecting click events) to produce
candidate click intervals with loose boundaries; Stage~2 applies a lightweight
boundary-refinement model --- for example, a shallow recurrent network or a
small attention block --- that takes a narrow context window around each candidate
boundary and emits a sub-frame corrected onset time. The Stage~2 model could be
trained on residuals from the Stage~1 predictions, focusing its capacity entirely
on the hard disambiguation cases near onset boundaries.

\paragraph{Hybrid detection-and-refinement strategy.}
Section~\ref{sec:eval-discussion} notes that the heuristic's boundary estimates,
when it does correctly match a click, are as tight as the RF's despite the
heuristic's far lower detection rate --- because the heuristic anchors its
duration estimate on the OS digital event, a reliable proxy that the learned
models do not use. This suggests a hybrid strategy that combines the strengths
of both: use the learned classifier to detect the clicks the OS event cannot
register (firmware-unregistered or noise-confounded presses), and fall back to
refining the boundary estimate against the OS digital event whenever it is
available, rather than relying purely on the analog depth signal for boundary
placement in cases the OS already resolves reliably.

\paragraph{FCN with systematic hyperparameter search.}
With a larger corpus (achievable by extending collection to additional hardware
platforms or broader user populations), a full Optuna sweep over the FCN's
architecture and training hyperparameters, run on the existing RTX~5080, would
determine whether the left-channel performance gap is closeable. The probe
ablation finding --- that
the RF's $\pm$10\,ms probe window captures most of the available local context ---
suggests that the FCN's 3\,s receptive field advantage is currently nullified by
data sparsity rather than by an architectural mismatch, making data expansion the
higher-leverage investment.

\paragraph{Active learning over cross-validation folds.}
The failure-mode analysis in Section~\ref{sec:eval-discussion} shows that one
held-out participant group (fold~4) drives most of the grouped cross-validation
variance for three of four tree-based model/channel combinations, consistent
with that group's profile diverging from the rest of the training distribution.
Rather than collecting additional data uniformly, an active-learning loop could
prioritize annotation effort toward participants whose telemetry is identified
as most dissimilar from the current training distribution, directly targeting
the atypical profiles responsible for the worst-case fold variance instead of
diluting the corpus with more typical examples.

\paragraph{Adequately powered, within-subject user study.}
The user study reported in Section~\ref{sec:eval-user-study} establishes a
significant CPS effect but was underpowered to resolve the Tracking and Flicking
results. A follow-up study should (1) recruit approximately 58 participants
(29 per group, per the Lehr's-formula estimate at the observed effect sizes),
(2) have each participant play every task twice, once under the placebo condition
and once under the calibrated condition, to reduce cross-participant skill-level
variance rather than relying purely on between-subject comparison, (3) introduce
a warmup round before each timed task to mitigate the learning-effect confound
between the calibration and test rounds, and (4) balance gamer and non-gamer
counts across groups in advance, or explicitly match participants on
self-reported skill level, insofar as recruitment constraints allow.

\paragraph{Movement-aware optimization objective.}
The Flicking result (Section~\ref{sec:eval-user-study}) suggests that optimizing
actuation timing in isolation may trade away accuracy in movement-heavy click
contexts: a lower actuation point shortens make latency but also shortens the
travel margin before an unintended actuation during a fast flick. Incorporating
mouse-movement telemetry (velocity or acceleration at the moment of click) into
the optimizer's objective function could let it distinguish stationary clicks,
where a low actuation point is unambiguously beneficial, from in-motion clicks,
where a more conservative threshold may better preserve accuracy.

\paragraph{Trajectory-aware and progressive sensitivity delivery.}
Beyond feeding movement telemetry into the objective function above, the same
signal could shape \emph{how} a recommended configuration is delivered rather
than only what is recommended. When several near-optimal (AP, RT) candidates
trade off similarly under the FSM simulator's loss, a trajectory-aware
recommender could favor whichever one least disrupts the player's existing
aim and flick habits. Separately, instead of applying a new configuration as
a single abrupt change, the system could ramp the actuation point toward the
recommended target progressively across sessions, giving the player's motor
system time to adapt and potentially avoiding the accuracy cost observed in
the Flicking task (Section~\ref{sec:eval-user-study}).

\section{Closing Remarks}

Returning to the hypothesis that opened this thesis (Chapter~\ref{chap:introduction}) ---
that per-user variation in click telemetry is structured enough to support
automatic labeling, statistical modeling, and parameter optimization beyond
static thresholds --- the evidence gathered here supports it, subject to the
qualifications catalogued above. The three contributions are not independent
results but stages of one pipeline: the labeling pipeline supplies the corpus
the classifier is trained on, the classifier's boundary predictions are what
the optimizer's FSM simulation consumes, and the user study is the first test
of whether these gains compound end-to-end rather than remaining theoretical.
That chain is validated most completely on the left channel, where labeling,
classification, and a statistically significant behavioral effect all align;
it is validated only in principle on the right channel and under cross-session
deployment, where data scarcity and untested robustness leave the thesis's
strongest numbers provisional rather than settled. The contribution of this
work is therefore best read as a validated methodology --- turning raw
per-user click telemetry into an actionable firmware recommendation --- together
with a precise map of where that methodology's evidence is strong and where it
still needs to be extended before the static-threshold status quo it targets
can be fully retired.
